\documentclass[12pt, a5paper]{article}

\usepackage{geometry}
\usepackage[utf8]{inputenc}
\usepackage{amssymb, amsmath, amsbsy}
\usepackage{mathpazo}
\usepackage{verbatim}

\renewcommand{\familydefault}{\sfdefault}

\begin{document}
% Luis Carrillo Gutiérrez
% \chapter*{Conjuntos}
\begin{center}
\subsection*{Conjuntos}
\end{center}

\paragraph{Objetivo 1 } -- Recordarás la definición de un conjunto y sus elementos.

\subsubsection*{Ejercicios resueltos :}

\begin{itemize}
\item {$\{2, 4, 6\}$ es un conjunto. Los elementos que forman este conjunto son : $2$, $4$, $6$.}

\item {¿Cuántos elementos hay en el conjunto $\{ manzana, pastel, durazno\}$? 3 elementos}

\item {$A = \{1, 2, 3 \}$ $B = \{2, 3, 4 \}$ ¿4 es un elemento de $A$? No
¿4 es un elemento de $B$? Sí }

\item {Si $U = \{ 1, 2, 3, 4, 5, 6 \}$, entonces $7 \notin U$,
¿Se podría extraer $A = \{ 1, 2, 3, 7 \}$ de este universo? No ¿Se podría extraer $B = \{ 2, 5 , 6 \}$? Sí }

\item{$A = \{5, 6, 7 \}$ $B = \{6, 7, 8 \}$  ¿$8 \in A$? No ¿$8 \in B$? Sí}

\item{Del ejemplo anterior como 8 no es un miembro de A podemos escribir:
$8 \notin A$}

\item{$A= \{ 1, 2, 3 \}$, $B = \{ 1, 5, 2, 7 \}$

% ¿Se cumple x ∈ A → x ∈ B ? SI
% ¿Se cumple x ∈ B → x ∈ A ? NO

¿Son iguales los dos conjuntos? NO
}
\item{$C = {6, 4}$ Escribe un conjunto D tal que D=C D = {4,6}}
\item{Si $U = \{ 1, 2, 3, 4, 5 \}$, $B = \{ 1, 2 \}$ y C = {3,4}, entonces el conjunto formado por todos los elementos comunes a B y C se le llama conjunto vació.}
\item{Si $P = \{ x | es un rio de la Tierra \}$, P también es finito aunque sea difícil contar los ríos del Mundo.}
\item{El conjunto de números que son múltiplos de 5 es un conjunto infinito porque no nunca se llega a un fin , observa: $A = \{ 5, 10, 15,  20, ...... \}$

}
\end{itemize}
\begin{comment}









Objetivo 2. Entenderás un conjunto de forma extensiva y comprensiva.

Ejercicios resueltos:1. Enunciar con palabras los siguientes incisos con el método de extensión
a) A = {x | x 2 = 4}
Se lee “A es el conjunto de los x tales que x al cuadrado es igual a cuatro”.
Los únicos números que elevados al cuadrado dan cuatro son 2 y -2, así que
A = {2, −2} .
b) B = {x | x − 2 = 5}
Se lee “B es el conjunto de los x tales que x menos 2 es igual a 5”. La única
solución es 7, de modo que B = {7} .
c) C = {x | x es positivo, x es negativo}
Se lee “C es le conjunto de los x tales que x es positivo y x es negativo”. No
hay ninguno número que sea positivo y negativo, así que C es vacío, es
decir, C = ∅ .
d) D = {x | x es una lera de la palabra "correcto"}
Se lee “D es el conjunto de los x tales que x es una letra de la palabra
correcto”. Las letras indicadas son c, o, r, e y t; así pues, D = {c, o, r, e, t} .

2. Escribir estos conjuntos con el método de compresión
a) A que consiste de las letras a, b, c, d y e. Pueden existir muchas soluciones 
primer resultado:
A = {x | x esta antes de f en el alfabeto} y como segundo resultado se tiene el
siguiente:
A = {x | x es unas de las primeras cinco letras del alfabeto}
b) B = {2, 4, 6,8,...}
B = {x | x es positivo y par}
c) El conjunto C de todos los países de Estados Unidos.
C = {x | x es un pais, x esta en los EstadosUnidos}
d) El conjunto D = {3}D = {x | x − 2 = 1} = {x | 2x = 6}
Objetivo 3. Recordaras la definición de subconjunto y la igualdad entre ellos.
Ejercicios resueltos:
1. Considere los siguientes conjuntos:
∅, A = {1} , B = {1,3} , C = {1,5,9} , D = {1, 2,3, 4,5} , E = {1,3,5, 7,9} , U = {1, 2,......,8,9}
Inserte el símbolo correcto ⊂ o ⊄ entre cada pareja de conjuntos:
(a ) ∅ ⊂ A ( b) A ⊂ B (c) B ⊄ C (d ) B ⊂ E
(e) C ⊄ D (f ) C ⊂ E (g ) D ⊄ E ( h ) D ⊂ U
a) ∅ ⊂ A ya que ∅ es un subconjunto de todo conjunto.
b) A ⊂ B ya que 1 es el único elemento de A y pertenece a B.
c) B ⊄ C ya que 3 ⊂ B pero 3 ∉ C .
d) B ⊂ E ya que los elementos de B también pertenecen a E.
e) C ⊄ D ya que 9 ∈ C pero 9 ∉ D .
f)
C ⊂ E ya que los elementos de C también pertenecen a E
g) D ⊄ E ya que 2 ∈ D pero 2 ∉ E .
h) D ⊂ U por que los elementos de D también pertenecen a U.
2. Considérese los conjuntos:
A = {1,3, 4,5,8,9} , B = {1, 2,3,5, 7} y C = {1,5}
Verificar si:
a) C ⊂ A y C ⊂ BSi se cumple ya que 1y 5 son elementos de A, B y C.
b) B ⊄ A
Si se cumple ya que 2 y 7, no pertenecen a A
Se puede observar que C ⊂ A
3. Usando los conjuntos dados, contesta si o no a las siguientes preguntas:
A = {1, 4, 2, 6,8,10} , B = {1, 4, 6,10} ,
C = {6, 4,1,10} , D = {6, 4,1}
U = {1, 2,3, 4,5, 6, 7,8,9,10}
¿Es A = D ?
¿Es D ⊆ A ?
¿Es B = C?
¿Es B ⊆ A?
¿Es A ⊆ B?
¿ Es A ≠ B?
¿ Es B ⊄ D?
¿ Es ∅ ⊆ D?
¿ Es ∅ = B?
¿ Es ∅ ⊆ B?
¿ Es B ⊆ U?
¿ Es A = U?
NO
SI
SI
SI
NO
SI
NO
SI
NO
SI
SI
NO

Objetivo 5. Recordarás las operaciones de unión e intersección de conjuntos y sus propiedades.

Ejemplos resueltos:

1. En el diagrama de Venn de la figura A ∪ B aparece rayado, o sea el área de A y el área de B.2. Sean S = {a, b, c, d} T = {f , b, d, g} . Entonces
y S ∪ T = {a, b, c, d, f , g}

3. Sea p el conjunto de los números reales positivos y Q el conjunto de los números
reales negativos P ∪ Q , unión de P y Q consiste en todos los números reales
exceptuando el cero.
La unión A y B se puede definir también concisamente así:
A ∪ B = {x |x ∈ A o x ∈ B}
Ejemplos resueltos
1. En el diagrama de Venn se ha rayado A ∩ B , el área común a ambos conjuntos
A y B.
2. Sea S = {a, b, c, d}
y
T = {f , b, d, g} . Entonces:
S ∩ T = {b, d}
3. Sean V = {2, 4, 6,...} , es decir los múltiplos de 2: y sea W = {3, 6,9,...} o sean
los múltiplos de 3. Entonces:V ∩ W = {6,12,18,...}
La intersección de A y B también se pueden definir concisamente así;
A ∩ B = {x |x ∈ A y x ∈ B}
Objetivo 6. Recordarás las operaciones diferencia y complemento de conjuntos y
sus propiedades.
Ejemplos resueltos:
1. En el diagrama de Venn se ha rayado el complemento de A, o sea el área exterior
de A. se supone que el conjunto universal U es el área del rectángulo.
2. Suponiendo que el conjunto universal U
sea el alfabeto, dado
T = {a, b, c} ,
entonces:
El complemento de T ' = {d, e, f , g, h,.....}
Ejemplos resueltos:
1. En el diagrama de Venn se ha rayado A-B, el área de A que nos es parte de B.2. Sean S = {a, b, c, d}
T = {f , b, d, g} . Se tiene:
y
S − T = {a, c}
3. Sea R el conjunto de los números reales y Q el conjunto de los números racionales.
Entonces R-Q es el conjunto de los números irracionales.
La diferencia de A y B se puede también definir concisamente como:
A − B = {x| x ∈ A y x ∉ B}
Objetivo 7. Recordarás la operación producto cruz, cardinalidad y potencia de
conjuntos y sus propiedades.
Ejemplos resueltos:
1. Determine el conjunto potencia P(S) de S = {a ,b ,c , d } los elementos de P(S) son
subconjuntos S. Así que:
P(S ) =
⎡ S ,{a ,b ,c},{a ,b , d },{a ,c , d },{b ,c , d },{a ,b},{a ,c},{a , d },{b ,c},{b , d },{c , d },{a},{b},{c},{d }, ∅ ⎤
⎣
⎦
Observa que P(S) tiene 2
4
=16 elementos.
s
{
}
2. Hallar el conjunto potencia 2 del conjunto S = 3,{1,4}Observar primero que S contiene dos elementos, 3 y el conjunto {1,4} . Por tanto,
s
2 contiene 2
2
= 4 elementos los cuales son:
{
}
2 = S ,{3},{{4}}, ∅
s
Ejemplos resueltos:
1. Sean W = { Juan, Josue, Ernesto} y V = {Maria, Carmen} . Hallar W x V.
W x V consiste en todos los pares ordenados (a, b) en los que a ∈ W y b ∈ V .
Por tanto:
{
}
W x V = (Juan, Maria), (Juan, Carmen), (Josue, Maria), (Josue, Carmen), (Ernesto, Maria), (Ernesto, Carmen)
2. Sean A = {a, b} , B = {2,3} ,
C = {3, 4} Hallar:
(a) A x (B ∪ C)
(b) (A x B) ∪ (A x C)
(c) A x (B ∩ C)
(d) (A x B) ∩ (A x C)
(a) A x (B ∪ C)
Se averigua primero B ∪ C = {2,3, 4} . Entonces:
A x (B ∪ C) = {(a, 2), (a,3), (a, 4), (b, 2), (b,3), (b, 4)}
(b) (A x B) ∪ (A x C)
Calcular primero A x B y A x C :
A x B = {(a, 2), (a,3), (b, 2), (b,3)}
A x C = {(a,3), (a, 4), (b,3), (b, 4)}
Ahora la unión de los conjuntos:
(A x B) ∪ (A x C) = {(a, 2), (a,3), (b, 2), (b,3), (a, 4), (b, 4)}
Los ejercicios (a) y (b) muestran que:A x (B ∪ C) = (A x B) ∪ (A x C)
(c) A x (B ∩ C)
Calcular primero B ∩ C = {3} .Entonces:
A x (B ∩ C) = {(a,3), (b,3)}
(d) (A x B) ∩ (A x C)
En (b) se calcularon A x B y A x C. La intersección de A x B y A x C es
el conjunto de los pares ordenados que pertenecen a ambos conjuntos, es
decir,
(A x B) ∩ (A x C) = {(a,3), (b,3)}
Por lo que (c) y (d) muestran que:
A x (B ∩ C) = (A x B) ∩ (A x C)
\end{comment}

\end{document}
